%%=============================================================================
%% LaTeX sjabloon voor bachelorproef, HoGent Bedrijf en Organisatie
%% Opleiding Toegepaste Informatica
%%=============================================================================

\documentclass[fleqn,a4paper,12pt]{book}

%%=============================================================================
%% LaTeX sjabloon voor de bachelorproef, HoGent Bedrijf en Organisatie
%% Opleiding toegepaste informatica
%%
%% Structuur en algemene vormgeving. Meestal hoef je hier niets te wijzigen.
%%
%% Vormgeving gebaseerd op "The Legrand Orange Book", version 2.0 (9/2/15)
%% door Mathias Legrand (legrand.mathias@gmail.com) met aanpassingen door
%% Vel (vel@latextemplates.com). Het oorspronkelijke template is te vinden op
%% http://www.LaTeXTemplates.com
%%
%% Aanpassingen voor HoGent toegepaste informatica: 
%%   Bert Van Vreckem <bert.vanvreckem@hogent.be>
%% Licentie: 
%%   CC BY-NC-SA 3.0 (http://creativecommons.org/licenses/by-nc-sa/3.0/)
%%=============================================================================

\nonstopmode

\usepackage{longtable}
\usepackage[acronym]{glossaries}
\usepackage{listings}  
\usepackage[final]{pdfpages}
\usepackage{scrextend}
\usepackage{listings} %For code in appendix
\lstset
{ %Formatting for code in appendix
    language=Matlab,
    basicstyle=\footnotesize,
    numbers=left,
    stepnumber=1,
    showstringspaces=false,
    tabsize=1,
    breaklines=true,
    breakatwhitespace=false,
}
\usepackage[acronym,xindy]{glossaries}

\makeglossaries

%%-----------------------------------------------------------------------------
%% Packages
%%-----------------------------------------------------------------------------

\usepackage[top=3cm,bottom=3cm,left=3cm,right=3cm,headsep=10pt,a4paper]{geometry} % Page margins
\usepackage[utf8]{inputenc}  % Accenten gebruiken in tekst (vb. é ipv \'e)
\usepackage{amsfonts}        % AMS math packages: extra wiskundige
\usepackage{amsmath}         %   symbolen (o.a. getallen-
\usepackage{amssymb}         %   verzamelingen N, R, Z, Q, etc.)
\usepackage[english,dutch]{babel}    % Taalinstellingen: woordsplitsingen,
                             %  commando's voor speciale karakters
                             %  ("dutch" voor NL)
\usepackage{iflang}
\usepackage{eurosym}         % Euro-symbool €
\usepackage{geometry}
\usepackage{graphicx}        % Invoegen van tekeningen
\graphicspath{{img/}}       % Specifies the directory where pictures are stored
\usepackage{tikz}            % Required for drawing custom shapes
\usepackage[pdftex,bookmarks=true]{hyperref}
                             % PDF krijgt klikbare links & verwijzingen,
                             %  inhoudstafel
\usepackage{enumitem}        % Customize lists
\setlist{nolistsep}         % Reduce spacing between list items
\usepackage{listings}        % Broncode mooi opmaken
\usepackage{multirow}        % Tekst over verschillende cellen in tabellen
\usepackage{rotating}        % Tabellen en figuren roteren

\usepackage{booktabs}        % Required for nicer horizontal rules in tables

\usepackage{xcolor}          % Required for specifying colors by name
\definecolor{maincolor}{RGB}{0,147,208} % Define the main color used for 
                             % highlighting throughout the book
                             % 0, 147, 208 = officiële kleur HoGent FBO

% Paragraph style: no indent, add space between paragraphs
\setlength{\parindent}{0em}
\setlength{\parskip}{1em}

\usepackage{etoolbox}
\usepackage{titling} % Macros for title, author, etc
\usepackage{lipsum}          % Voor vultekst (lorem ipsum)
\usepackage{underscore} % make it possible to use _ in your text

%----------------------------------------------------------------------------------------
%	FONTS
%----------------------------------------------------------------------------------------

\usepackage{avant} % Use the Avantgarde font for headings
%\usepackage{times} % Use the Times font for headings
\usepackage{mathptmx} % Use the Adobe Times Roman as the default text font together with math symbols from the Sym­bol, Chancery and Com­puter Modern fonts

\usepackage{microtype} % Slightly tweak font spacing for aesthetics
\usepackage[utf8]{inputenc} % Required for including letters with accents
\usepackage[T1]{fontenc} % Use 8-bit encoding that has 256 glyphs

%------------------------------------------------------------------------------
%	TITLE PAGE
%------------------------------------------------------------------------------

\newcommand{\inserttitlepage}{%
\begin{titlepage}
  \newgeometry{top=2cm,bottom=1.5cm,left=1.5cm,right=1.5cm}
  \begin{center}

    \begingroup
    \rmfamily
    \includegraphics[width=2.5cm]{img/HG-beeldmerk-woordmerk}\\[.5cm]
    Faculteit Bedrijf en Organisatie\\[3cm]
    \titel
    \vfill
    \student\\[3.5cm]
    Scriptie voorgedragen tot het bekomen van de graad van\\professionele bachelor in de toegepaste informatica\\[2cm]
    Promotor:\\
    \promotor\\
    \ifdefempty{\copromotor}{\vspace{2.5cm}}{Co-promotor:\\\copromotor\\[2.5cm]}
    Instelling: \instelling\\[.5cm]
    Academiejaar: \academiejaar\\[.5cm]
    \ifcase \examenperiode \or Eerste \or Tweede \else Derde \fi examenperiode
    \endgroup

  \end{center}
  \restoregeometry
\end{titlepage}
  \emptypage
\begin{titlepage}
  \newgeometry{top=5.35cm,bottom=1.5cm,left=1.5cm,right=1.5cm}
  \begin{center}

    \begingroup
    \rmfamily
    \IfLanguageName{dutch}{Faculteit Bedrijf en Organisatie}{Faculty of Business and Information Management}\\[3cm]
    \titel
    \vfill
    \student\\[3.5cm]
    \IfLanguageName{dutch}{Scriptie voorgedragen tot het bekomen van de graad van\\professionele bachelor in de toegepaste informatica}{Thesis submitted in partial fulfillment of the requirements for the degree of\\professional bachelor of applied computer science}\\[2cm]
    Promotor:\\
    \promotor\\
    \ifdefempty{\copromotor}{\vspace{2.5cm}}{Co-promotor:\\\copromotor\\[2.5cm]}
    \IfLanguageName{dutch}{Instelling}{Institution}: \instelling\\[.5cm]
    \IfLanguageName{dutch}{Academiejaar}{Academic year}: \academiejaar\\[.5cm]
    \IfLanguageName{dutch}{%
    \ifcase \examenperiode \or Eerste \or Tweede \else Derde \fi examenperiode}{%
    \ifcase \examenperiode \or First \or Second \else Third \fi examination period}
    \endgroup

  \end{center}
  \restoregeometry
\end{titlepage}
}

%----------------------------------------------------------------------------------------
%	BIBLIOGRAPHY AND INDEX
%----------------------------------------------------------------------------------------

\usepackage[style=apa,backend=biber]{biblatex}
\usepackage{csquotes}
\DeclareLanguageMapping{dutch}{dutch-apa}
\addbibresource{bachproef-tin.bib} % BibTeX bibliography file
\defbibheading{bibempty}{}

\usepackage{calc} % For simpler calculation - used for spacing the index letter headings correctly
\usepackage{makeidx} % Required to make an index
\makeindex % Tells LaTeX to create the files required for indexing

%----------------------------------------------------------------------------------------
%	MAIN TABLE OF CONTENTS
%----------------------------------------------------------------------------------------

\usepackage{titletoc} % Required for manipulating the table of contents

\contentsmargin{0cm} % Removes the default margin

% Part text styling
\titlecontents{part}[0cm]
{\addvspace{20pt}\centering\large\bfseries}
{}
{}
{}

% Chapter text styling
\titlecontents{chapter}[1.25cm] % Indentation
{\addvspace{12pt}\large\sffamily\bfseries} % Spacing and font options for chapters
{\color{maincolor!60}\contentslabel[\Large\thecontentslabel]{1.25cm}\color{maincolor}} % Chapter number
{\color{maincolor}}
{\color{maincolor!60}\normalsize\;\titlerule*[.5pc]{.}\;\thecontentspage} % Page number

% Section text styling
\titlecontents{section}[1.25cm] % Indentation
{\addvspace{3pt}\sffamily\bfseries} % Spacing and font options for sections
{\contentslabel[\thecontentslabel]{1.25cm}} % Section number
{}
{\hfill\color{black}\thecontentspage} % Page number
[]

% Subsection text styling
\titlecontents{subsection}[1.25cm] % Indentation
{\addvspace{1pt}\sffamily\small} % Spacing and font options for subsections
{\contentslabel[\thecontentslabel]{1.25cm}} % Subsection number
{}
{\ \titlerule*[.5pc]{.}\;\thecontentspage} % Page number
[]

% List of figures
\titlecontents{figure}[0em]
{\addvspace{-5pt}\sffamily}
{\thecontentslabel\hspace*{1em}}
{}
{\ \titlerule*[.5pc]{.}\;\thecontentspage}
[]

% List of tables
\titlecontents{table}[0em]
{\addvspace{-5pt}\sffamily}
{\thecontentslabel\hspace*{1em}}
{}
{\ \titlerule*[.5pc]{.}\;\thecontentspage}
[]

%----------------------------------------------------------------------------------------
%	MINI TABLE OF CONTENTS IN PART HEADS
%----------------------------------------------------------------------------------------

% Chapter text styling
\titlecontents{lchapter}[0em] % Indenting
{\addvspace{15pt}\large\sffamily\bfseries} % Spacing and font options for chapters
{\color{maincolor}\contentslabel[\Large\thecontentslabel]{1.25cm}\color{maincolor}} % Chapter number
{}
{\color{maincolor}\normalsize\sffamily\bfseries\;\titlerule*[.5pc]{.}\;\thecontentspage} % Page number

% Section text styling
\titlecontents{lsection}[0em] % Indenting
{\sffamily\small} % Spacing and font options for sections
{\contentslabel[\thecontentslabel]{1.25cm}} % Section number
{}
{}

% Subsection text styling
\titlecontents{lsubsection}[.5em] % Indentation
{\normalfont\footnotesize\sffamily} % Font settings
{}
{}
{}

%----------------------------------------------------------------------------------------
%	PAGE HEADERS
%----------------------------------------------------------------------------------------

\usepackage{fancyhdr} % Required for header and footer configuration

\pagestyle{fancy}
\renewcommand{\chaptermark}[1]{\markboth{\sffamily\normalsize\bfseries\chaptername\ \thechapter.\ #1}{}} % Chapter text font settings
\renewcommand{\sectionmark}[1]{\markright{\sffamily\normalsize\thesection\hspace{5pt}#1}{}} % Section text font settings
\fancyhf{} \fancyhead[LE,RO]{\sffamily\normalsize\thepage} % Font setting for the page number in the header
\fancyhead[LO]{\rightmark} % Print the nearest section name on the left side of odd pages
\fancyhead[RE]{\leftmark} % Print the current chapter name on the right side of even pages
\renewcommand{\headrulewidth}{0.5pt} % Width of the rule under the header
\addtolength{\headheight}{2.5pt} % Increase the spacing around the header slightly
\renewcommand{\footrulewidth}{0pt} % Removes the rule in the footer
\fancypagestyle{plain}{\fancyhead{}\renewcommand{\headrulewidth}{0pt}} % Style for when a plain pagestyle is specified

% Removes the header from odd empty pages at the end of chapters
\makeatletter
\renewcommand{\cleardoublepage}{
\clearpage\ifodd\c@page\else
\hbox{}
\vspace*{\fill}
\thispagestyle{empty}
\newpage
\fi}

%----------------------------------------------------------------------------------------
%	THEOREM STYLES
%----------------------------------------------------------------------------------------

\usepackage{amsmath,amsfonts,amssymb,amsthm} % For math equations, theorems, symbols, etc

\newcommand{\intoo}[2]{\mathopen{]}#1\,;#2\mathclose{[}}
\newcommand{\ud}{\mathop{\mathrm{{}d}}\mathopen{}}
\newcommand{\intff}[2]{\mathopen{[}#1\,;#2\mathclose{]}}
\newtheorem{notation}{Notation}[chapter]

% Boxed/framed environments
\newtheoremstyle{maincolornumbox}% % Theorem style name
{0pt}% Space above
{0pt}% Space below
{\normalfont}% % Body font
{}% Indent amount
{\small\bf\sffamily\color{maincolor}}% % Theorem head font
{\;}% Punctuation after theorem head
{0.25em}% Space after theorem head
{\small\sffamily\color{maincolor}\thmname{#1}\nobreakspace\thmnumber{\@ifnotempty{#1}{}\@upn{#2}}% Theorem text (e.g. Theorem 2.1)
\thmnote{\nobreakspace\the\thm@notefont\sffamily\bfseries\color{black}---\nobreakspace#3.}} % Optional theorem note
\renewcommand{\qedsymbol}{$\blacksquare$}% Optional qed square

\newtheoremstyle{blacknumex}% Theorem style name
{5pt}% Space above
{5pt}% Space below
{\normalfont}% Body font
{} % Indent amount
{\small\bf\sffamily}% Theorem head font
{\;}% Punctuation after theorem head
{0.25em}% Space after theorem head
{\small\sffamily{\tiny\ensuremath{\blacksquare}}\nobreakspace\thmname{#1}\nobreakspace\thmnumber{\@ifnotempty{#1}{}\@upn{#2}}% Theorem text (e.g. Theorem 2.1)
\thmnote{\nobreakspace\the\thm@notefont\sffamily\bfseries---\nobreakspace#3.}}% Optional theorem note

\newtheoremstyle{blacknumbox} % Theorem style name
{0pt}% Space above
{0pt}% Space below
{\normalfont}% Body font
{}% Indent amount
{\small\bf\sffamily}% Theorem head font
{\;}% Punctuation after theorem head
{0.25em}% Space after theorem head
{\small\sffamily\thmname{#1}\nobreakspace\thmnumber{\@ifnotempty{#1}{}\@upn{#2}}% Theorem text (e.g. Theorem 2.1)
\thmnote{\nobreakspace\the\thm@notefont\sffamily\bfseries---\nobreakspace#3.}}% Optional theorem note

% Non-boxed/non-framed environments
\newtheoremstyle{maincolornum}% % Theorem style name
{5pt}% Space above
{5pt}% Space below
{\normalfont}% % Body font
{}% Indent amount
{\small\bf\sffamily\color{maincolor}}% % Theorem head font
{\;}% Punctuation after theorem head
{0.25em}% Space after theorem head
{\small\sffamily\color{maincolor}\thmname{#1}\nobreakspace\thmnumber{\@ifnotempty{#1}{}\@upn{#2}}% Theorem text (e.g. Theorem 2.1)
\thmnote{\nobreakspace\the\thm@notefont\sffamily\bfseries\color{black}---\nobreakspace#3.}} % Optional theorem note
\renewcommand{\qedsymbol}{$\blacksquare$}% Optional qed square
\makeatother

% Defines the theorem text style for each type of theorem to one of the three styles above
\newcounter{dummy}
\numberwithin{dummy}{section}
\theoremstyle{maincolornumbox}
\newtheorem{theoremeT}[dummy]{Theorem}
\newtheorem{problem}{Problem}[chapter]
\newtheorem{exerciseT}{Exercise}[chapter]
\theoremstyle{blacknumex}
\newtheorem{exampleT}{Example}[chapter]
\theoremstyle{blacknumbox}
\newtheorem{vocabulary}{Vocabulary}[chapter]
\newtheorem{definitionT}{Definition}[section]
\newtheorem{corollaryT}[dummy]{Corollary}
\theoremstyle{maincolornum}
\newtheorem{proposition}[dummy]{Proposition}

%----------------------------------------------------------------------------------------
%	DEFINITION OF COLORED BOXES
%----------------------------------------------------------------------------------------

\RequirePackage[framemethod=default]{mdframed} % Required for creating the theorem, definition, exercise and corollary boxes

% Theorem box
\newmdenv[skipabove=7pt,
skipbelow=7pt,
backgroundcolor=black!5,
linecolor=maincolor,
innerleftmargin=5pt,
innerrightmargin=5pt,
innertopmargin=5pt,
leftmargin=0cm,
rightmargin=0cm,
innerbottommargin=5pt]{tBox}

% Exercise box
\newmdenv[skipabove=7pt,
skipbelow=7pt,
rightline=false,
leftline=true,
topline=false,
bottomline=false,
backgroundcolor=maincolor!10,
linecolor=maincolor,
innerleftmargin=5pt,
innerrightmargin=5pt,
innertopmargin=5pt,
innerbottommargin=5pt,
leftmargin=0cm,
rightmargin=0cm,
linewidth=4pt]{eBox}

% Definition box
\newmdenv[skipabove=7pt,
skipbelow=7pt,
rightline=false,
leftline=true,
topline=false,
bottomline=false,
linecolor=maincolor,
innerleftmargin=5pt,
innerrightmargin=5pt,
innertopmargin=0pt,
leftmargin=0cm,
rightmargin=0cm,
linewidth=4pt,
innerbottommargin=0pt]{dBox}

% Corollary box
\newmdenv[skipabove=7pt,
skipbelow=7pt,
rightline=false,
leftline=true,
topline=false,
bottomline=false,
linecolor=gray,
backgroundcolor=black!5,
innerleftmargin=5pt,
innerrightmargin=5pt,
innertopmargin=5pt,
leftmargin=0cm,
rightmargin=0cm,
linewidth=4pt,
innerbottommargin=5pt]{cBox}

% Creates an environment for each type of theorem and assigns it a theorem text style from the "Theorem Styles" section above and a colored box from above
\newenvironment{theorem}{\begin{tBox}\begin{theoremeT}}{\end{theoremeT}\end{tBox}}
\newenvironment{exercise}{\begin{eBox}\begin{exerciseT}}{\hfill{\color{maincolor}\tiny\ensuremath{\blacksquare}}\end{exerciseT}\end{eBox}}
\newenvironment{definition}{\begin{dBox}\begin{definitionT}}{\end{definitionT}\end{dBox}}
\newenvironment{example}{\begin{exampleT}}{\hfill{\tiny\ensuremath{\blacksquare}}\end{exampleT}}
\newenvironment{corollary}{\begin{cBox}\begin{corollaryT}}{\end{corollaryT}\end{cBox}}

%----------------------------------------------------------------------------------------
%	REMARK ENVIRONMENT
%----------------------------------------------------------------------------------------

\newenvironment{remark}{\par\vspace{10pt}\small % Vertical white space above the remark and smaller font size
\begin{list}{}{
\leftmargin=35pt % Indentation on the left
\rightmargin=25pt}\item\ignorespaces % Indentation on the right
\makebox[-2.5pt]{\begin{tikzpicture}[overlay]
\node[draw=maincolor!60,line width=1pt,circle,fill=maincolor!25,font=\sffamily\bfseries,inner sep=2pt,outer sep=0pt] at (-15pt,0pt){\textcolor{maincolor}{R}};\end{tikzpicture}} % Orange R in a circle
\advance\baselineskip -1pt}{\end{list}\vskip5pt} % Tighter line spacing and white space after remark

%----------------------------------------------------------------------------------------
%	SECTION NUMBERING IN THE MARGIN
%----------------------------------------------------------------------------------------

\makeatletter
\renewcommand{\@seccntformat}[1]{\llap{\textcolor{maincolor}{\csname the#1\endcsname}\hspace{1em}}}
\renewcommand{\section}{\@startsection{section}{1}{\z@}
{-4ex \@plus -1ex \@minus -.4ex}
{1ex \@plus.2ex }
{\normalfont\large\sffamily\bfseries}}
\renewcommand{\subsection}{\@startsection {subsection}{2}{\z@}
{-3ex \@plus -0.1ex \@minus -.4ex}
{0.5ex \@plus.2ex }
{\normalfont\sffamily\bfseries}}
\renewcommand{\subsubsection}{\@startsection {subsubsection}{3}{\z@}
{-2ex \@plus -0.1ex \@minus -.2ex}
{.2ex \@plus.2ex }
{\normalfont\small\sffamily\bfseries}}
\renewcommand\paragraph{\@startsection{paragraph}{4}{\z@}
{-2ex \@plus-.2ex \@minus .2ex}
{.1ex}
{\normalfont\small\sffamily\bfseries}}

%----------------------------------------------------------------------------------------
%	PART HEADINGS
%----------------------------------------------------------------------------------------

% numbered part in the table of contents
\newcommand{\@mypartnumtocformat}[2]{%
\setlength\fboxsep{0pt}%
\noindent\colorbox{maincolor!20}{\strut\parbox[c][.7cm]{\ecart}{\color{maincolor!70}\Large\sffamily\bfseries\centering#1}}\hskip\esp\colorbox{maincolor!40}{\strut\parbox[c][.7cm]{\linewidth-\ecart-\esp}{\Large\sffamily\centering#2}}}%
%%%%%%%%%%%%%%%%%%%%%%%%%%%%%%%%%%
% unnumbered part in the table of contents
\newcommand{\@myparttocformat}[1]{%
\setlength\fboxsep{0pt}%
\noindent\colorbox{maincolor!40}{\strut\parbox[c][.7cm]{\linewidth}{\Large\sffamily\centering#1}}}%
%%%%%%%%%%%%%%%%%%%%%%%%%%%%%%%%%%
\newlength\esp
\setlength\esp{4pt}
\newlength\ecart
\setlength\ecart{1.2cm-\esp}
\newcommand{\thepartimage}{}%
\newcommand{\partimage}[1]{\renewcommand{\thepartimage}{#1}}%
\def\@part[#1]#2{%
\ifnum \c@secnumdepth >-2\relax%
\refstepcounter{part}%
\addcontentsline{toc}{part}{\texorpdfstring{\protect\@mypartnumtocformat{\thepart}{#1}}{\partname~\thepart\ ---\ #1}}
\else%
\addcontentsline{toc}{part}{\texorpdfstring{\protect\@myparttocformat{#1}}{#1}}%
\fi%
\startcontents%
\markboth{}{}%
{\thispagestyle{empty}%
\begin{tikzpicture}[remember picture,overlay]%
\node at (current page.north west){\begin{tikzpicture}[remember picture,overlay]%
\fill[maincolor!20](0cm,0cm) rectangle (\paperwidth,-\paperheight);
\node[anchor=north] at (4cm,-3.25cm){\color{maincolor!40}\fontsize{220}{100}\sffamily\bfseries\@Roman\c@part};
\node[anchor=south east] at (\paperwidth-1cm,-\paperheight+1cm){\parbox[t][][t]{8.5cm}{
\printcontents{l}{0}{\setcounter{tocdepth}{1}}%
}};
\node[anchor=north east] at (\paperwidth-1.5cm,-3.25cm){\parbox[t][][t]{15cm}{\strut\raggedleft\color{white}\fontsize{30}{30}\sffamily\bfseries#2}};
\end{tikzpicture}};
\end{tikzpicture}}%
\@endpart}
\def\@spart#1{%
\startcontents%
\phantomsection
{\thispagestyle{empty}%
\begin{tikzpicture}[remember picture,overlay]%
\node at (current page.north west){\begin{tikzpicture}[remember picture,overlay]%
\fill[maincolor!20](0cm,0cm) rectangle (\paperwidth,-\paperheight);
\node[anchor=north east] at (\paperwidth-1.5cm,-3.25cm){\parbox[t][][t]{15cm}{\strut\raggedleft\color{white}\fontsize{30}{30}\sffamily\bfseries#1}};
\end{tikzpicture}};
\end{tikzpicture}}
\addcontentsline{toc}{part}{\texorpdfstring{%
\setlength\fboxsep{0pt}%
\noindent\protect\colorbox{maincolor!40}{\strut\protect\parbox[c][.7cm]{\linewidth}{\Large\sffamily\protect\centering #1\quad\mbox{}}}}{#1}}%
\@endpart}
\def\@endpart{\vfil\newpage
\if@twoside
\if@openright
\null
\thispagestyle{empty}%
\newpage
\fi
\fi
\if@tempswa
\twocolumn
\fi}

%----------------------------------------------------------------------------------------
%	CHAPTER HEADINGS
%----------------------------------------------------------------------------------------

% A switch to conditionally include a picture, implemented by  Christian Hupfer
\newif\ifusechapterimage
\usechapterimagetrue
\newcommand{\thechapterimage}{}%
\newcommand{\chapterimage}[1]{\ifusechapterimage\renewcommand{\thechapterimage}{#1}\fi}%
\def\@makechapterhead#1{%
{\parindent \z@ \raggedright \normalfont
\ifnum \c@secnumdepth >\m@ne
\if@mainmatter
\begin{tikzpicture}[remember picture,overlay]
\node at (current page.north west)
{\begin{tikzpicture}[remember picture,overlay]
\node[anchor=north west,inner sep=0pt] at (0,0) {\ifusechapterimage\includegraphics[width=\paperwidth]{\thechapterimage}\fi};
\draw[anchor=west] (\Gm@lmargin,-9cm) node [line width=2pt,rounded corners=15pt,draw=maincolor,fill=white,fill opacity=0.5,inner sep=15pt]{\strut\makebox[22cm]{}};
\draw[anchor=west] (\Gm@lmargin+.3cm,-9cm) node {\huge\sffamily\bfseries\color{black}\thechapter. #1\strut};
\end{tikzpicture}};
\end{tikzpicture}
\else
\begin{tikzpicture}[remember picture,overlay]
\node at (current page.north west)
{\begin{tikzpicture}[remember picture,overlay]
\node[anchor=north west,inner sep=0pt] at (0,0) {\ifusechapterimage\includegraphics[width=\paperwidth]{\thechapterimage}\fi};
\draw[anchor=west] (\Gm@lmargin,-9cm) node [line width=2pt,rounded corners=15pt,draw=maincolor,fill=white,fill opacity=0.5,inner sep=15pt]{\strut\makebox[22cm]{}};
\draw[anchor=west] (\Gm@lmargin+.3cm,-9cm) node {\huge\sffamily\bfseries\color{black}#1\strut};
\end{tikzpicture}};
\end{tikzpicture}
\fi\fi\par\vspace*{270\p@}}}

%-------------------------------------------

\def\@makeschapterhead#1{%
\begin{tikzpicture}[remember picture,overlay]
\node at (current page.north west)
{\begin{tikzpicture}[remember picture,overlay]
\node[anchor=north west,inner sep=0pt] at (0,0) {\ifusechapterimage\includegraphics[width=\paperwidth]{\thechapterimage}\fi};
\draw[anchor=west] (\Gm@lmargin,-9cm) node [line width=2pt,rounded corners=15pt,draw=maincolor,fill=white,fill opacity=0.5,inner sep=15pt]{\strut\makebox[22cm]{}};
\draw[anchor=west] (\Gm@lmargin+.3cm,-9cm) node {\huge\sffamily\bfseries\color{black}#1\strut};
\end{tikzpicture}};
\end{tikzpicture}
\par\vspace*{270\p@}}
\makeatother

%----------------------------------------------------------------------------------------
%	HYPERLINKS IN THE DOCUMENTS
%----------------------------------------------------------------------------------------

\usepackage{hyperref}
\hypersetup{hidelinks,backref=true,pagebackref=true,hyperindex=true,colorlinks=false,breaklinks=true,urlcolor= maincolor,bookmarks=true,bookmarksopen=false,pdftitle={Title},pdfauthor={Author}}
\usepackage{bookmark}
\bookmarksetup{
open,
numbered,
addtohook={%
\ifnum\bookmarkget{level}=0 % chapter
\bookmarksetup{bold}%
\fi
\ifnum\bookmarkget{level}=-1 % part
\bookmarksetup{color=maincolor,bold}%
\fi
}
}

%----------------------------------------------------------------------------------------
%	Java source code
%----------------------------------------------------------------------------------------

% Commando voor invoegen Java-broncodebestanden (dank aan Niels Corneille)
% Gebruik:
%   \codefragment{source/MijnKlasse.java}{Uitleg bij de code}
%
% Je kan dit aanpassen aan de taal die je zelf het meeste gebruikt in je
% bachelorproef.
\newcommand{\codefragment}[2]{ \lstset{%
  language=java,
  breaklines=true,
  float=th,
  caption={#2},
  basicstyle=\scriptsize,
  frame=single,
  extendedchars=\true
}
\lstinputlisting{#1}}

% Leeg blad
\newcommand{\emptypage}{%
\newpage
\thispagestyle{empty}
\mbox{}
\newpage
}


%%---------- Documenteigenschappen --------------------------------------------
%% TODO: Vul dit aan met je eigen info:

% Je eigen naam
\newcommand{\student}{Robin Malfait}

% De naam van je promotor (lector van de opleiding)
\newcommand{\promotor}{Lieven Smits}

% De naam van je co-promotor. Als je promotor ook je opdrachtgever is en je
% dus ook inhoudelijk begeleidt (en enkel dan!), mag je dit leeg laten.
\newcommand{\copromotor}{Mathias Verraes}

% Indien je bachelorproef in opdracht van/in samenwerking met een bedrijf of
% externe organisatie geschreven is, geef je hier de naam. Zoniet laat je dit
% zoals het is.
\newcommand{\instelling}{---}

% De titel van het rapport/bachelorproef
\newcommand{\titel}{Wanneer is EventSourcing een meerwaarde voor een bedrijf zoals Skedify, waar ze gespecialiseerd zijn in online scheduling?}

% Datum van indienen (gebruik telkens de deadline, ook al geef je eerder af)
\newcommand{\datum}{02 juni 2017}

% Academiejaar
\newcommand{\academiejaar}{2016-2017}

% Examenperiode
%  - 1e semester = 1e examenperiode => 1
%  - 2e semester = 2e examenperiode => 2
%  - tweede zit  = 3e examenperiode => 3
\newcommand{\examenperiode}{2}

%%=============================================================================
%% Inhoud document
%%=============================================================================

\begin{document}

%---------- Taalselectie ------------------------------------------------------
%% Als je je bachelorproef in het Engels schrijft, haal dan onderstaande regel
%% uit commentaar. Let op: de tekst op de voorkaft blijft in het Nederlands, en
%% dat is ook de bedoeling!
%\selectlanguage{english}

%---------- Titelblad ---------------------------------------------------------
\inserttitlepage

%---------- Samenvatting, voorwoord -------------------------------------------
\usechapterimagefalse
%%=============================================================================
%% Samenvatting
%%=============================================================================

%% TODO: De "abstract" of samenvatting is een kernachtige (~ 1 blz. voor een
%% thesis) synthese van het document.
%%
%% Deze aspecten moeten zeker aan bod komen:
%% - Context: waarom is dit werk belangrijk?
%% - Nood: waarom moest dit onderzocht worden?
%% - Taak: wat heb je precies gedaan?
%% - Object: wat staat in dit document geschreven?
%% - Resultaat: wat was het resultaat?
%% - Conclusie: wat is/zijn de belangrijkste conclusie(s)?
%% - Perspectief: blijven er nog vragen open die in de toekomst nog kunnen
%%    onderzocht worden? Wat is een mogelijk vervolg voor jouw onderzoek?
%%
%% LET OP! Een samenvatting is GEEN voorwoord!

%%---------- Nederlandse samenvatting -----------------------------------------
%%
%% TODO: Als je je bachelorproef in het Engels schrijft, moet je eerst een
%% Nederlandse samenvatting invoegen. Haal daarvoor onderstaande code uit
%% commentaar.
%% Wie zijn bachelorproef in het Nederlands schrijft, kan dit negeren en heel
%% deze sectie verwijderen.

\IfLanguageName{english}{%
\selectlanguage{dutch}
\chapter*{Samenvatting}
\lipsum[1-4]
\selectlanguage{english}
}{}

%%---------- Samenvatting -----------------------------------------------------
%%
%% De samenvatting in de hoofdtaal van het document

\chapter*{\IfLanguageName{dutch}{Samenvatting}{Abstract}}

\lipsum[1-4]

%%=============================================================================
%% Voorwoord
%%=============================================================================

\chapter*{Voorwoord}
\label{ch:voorwoord}

%% TODO:
%% Het voorwoord is het enige deel van de bachelorproef waar je vanuit je
%% eigen standpunt (``ik-vorm'') mag schrijven. Je kan hier bv. motiveren
%% waarom jij het onderwerp wil bespreken.
%% Vergeet ook niet te bedanken wie je geholpen/gesteund/... heeft



%---------- Inhoudstafel ------------------------------------------------------
\pagestyle{empty} % No headers
\tableofcontents % Print the table of contents itself
\cleardoublepage % Forces the first chapter to start on an odd page so it's on the right
\pagestyle{fancy} % Print headers again

%---------- Lijst afkortingen, termen -----------------------------------------
%% Als je een lijst van afkortingen of termen wil toevoegen, dan hoort die
%% hier thuis. Gebruik bijvoorbeeld de ``glossaries'' package.

%%---------- Kern -------------------------------------------------------------

%%=============================================================================
%% Inleiding
%%=============================================================================

\chapter{Inleiding}
\label{ch:inleiding}

De inleiding moet de lezer alle nodige informatie verschaffen om het onderwerp te begrijpen zonder nog externe werken te moeten raadplegen \autocite{Pollefliet2011}. Dit is een doorlopende tekst die gebaseerd is op al wat je over het onderwerp gelezen hebt (literatuuronderzoek).

Je verwijst bij elke bewering die je doet, vakterm die je introduceert, enz. naar je bronnen. In \LaTeX{} kan dat met het commando \texttt{$\backslash${textcite\{\}}} of \texttt{$\backslash${autocite\{\}}}. Als argument van het commando geef je de ``sleutel'' van een ``record'' in een bibliografische databank in het Bib\TeX{}-formaat (een tekstbestand). Als je expliciet naar de auteur verwijst in de zin, gebruik je \texttt{$\backslash${}textcite\{\}}.
Soms wil je de auteur niet expliciet vernoemen, dan gebruik je \texttt{$\backslash${}autocite\{\}}. Hieronder een voorbeeld van elk.

\textcite{Knuth1998} schreef een van de standaardwerken over sorteer- en zoekalgoritmen. Experten zijn het erover eens dat cloud computing een interessante opportuniteit vormen, zowel voor gebruikers als voor dienstverleners op vlak van informatietechnologie~\autocite{Creeger2009}.

\section{Stand van zaken}
\label{sec:stand-van-zaken}

%% TODO: deze sectie (die je kan opsplitsen in verschillende secties) bevat je
%% literatuurstudie. Vergeet niet telkens je bronnen te vermelden!

\section{Probleemstelling en Onderzoeksvragen}
\label{sec:onderzoeksvragen}

%% Uit je probleemstelling moet duidelijk zijn dat je onderzoek een meerwaarde
%% heeft voor een concrete doelgroep (bv. een bedrijf).
%%
%% Wees zo concreet mogelijk bij het formuleren van je
%% onderzoeksvra(a)g(en). Een onderzoeksvraag is trouwens iets waar nog
%% niemand op dit moment een antwoord heeft (voor zover je kan nagaan).
Data is belangrijk, rapportering is belangrijk. Voorspellen van vragen die de business kan vragen binnen dit en 5 jaar is onmogelijk. Daarom is het belangrijk om alle informatie bij te houden om hier interessante rapportering op te kunnen uitvoeren. In een systeem met een standaard relationele databank wordt niet alle informatie bijgehouden. Bij het verwijderen of wijzigen van informatie is de voorgaande informatie verloreren gegaan. Alle informatie bijhouden is dus interessant, maar wat zijn hier de voor en nadelen van?

De onderzoeksvraag luidt als volgt: Wanneer is EventSourcing een meerwaarde voor een bedrijf zoals Skedify, waar ze gespecialiseerd zijn in online scheduling?

Hierbij worden ook volgende subvragen beantwoord:

\begin{itemize}
    \item Wat is het verschil tussen een systeem met relationele databank modellen en een systeem met EventSourcing?
    \item Hoe kan een applicatie die gebruik maakt van EventSourcing getest worden?
    \item Welke delen moeten er EventSourced worden, en hoe kan je deze delen bepalen?
\end{itemize}

\section{Opzet van deze bachelorproef}
\label{sec:opzet-bachelorproef}

%% TODO: Het is gebruikelijk aan het einde van de inleiding een overzicht te
%% geven van de opbouw van de rest van de tekst. Deze sectie bevat al een aanzet
%% die je kan aanvullen/aanpassen in functie van je eigen tekst.

De rest van deze bachelorproef is als volgt opgebouwd:

In Hoofdstuk~\ref{ch:methodologie} wordt de methodologie toegelicht en worden de gebruikte onderzoekstechnieken besproken om een antwoord te kunnen formuleren op de onderzoeksvragen.

%% TODO: Vul hier aan voor je eigen hoofstukken, één of twee zinnen per hoofdstuk

In Hoofdstuk~\ref{ch:conclusie}, tenslotte, wordt de conclusie gegeven en een antwoord geformuleerd op de onderzoeksvragen. Daarbij wordt ook een aanzet gegeven voor toekomstig onderzoek binnen dit domein.


%%=============================================================================
%% Methodologie
%%=============================================================================

\chapter{Methodologie}
\label{ch:methodologie}

%% TODO: Hoe ben je te werk gegaan? Verdeel je onderzoek in grote fasen, en
%% licht in elke fase toe welke stappen je gevolgd hebt. Verantwoord waarom je
%% op deze manier te werk gegaan bent. Je moet kunnen aantonen dat je de best
%% mogelijke manier toegepast hebt om een antwoord te vinden op de
%% onderzoeksvraag.


Omdat er met EventSourcing gewerkt zal worden om een proof of concept te maken, moet er eerst gekeken worden naar wat EventSourcing is en wat het allemaal inhoudt. Naast EventSourcing zijn er ook concepten zoals CQRS die aan bod moeten komen. Daarnaast draait het in dit onderzoek rond Skedify, het bedrijf die als concrete case gebruikt zal worden. In het deel rond Skedify, kan er meer te weten gekomen worden in verband met wat ze doen en hoe ze te werk gaan. Daarna zal er gekeken worden naar de huidige manier van werken. Tot slot gaan we alles samen voegen: Skedify en de huidige manier van werken, en Skedify en EventSourcing. Hierna kunnen we dan de voor- en nadelen van beide mogelijkheden afgaan om te zien of EventSourcing nu effectief een meerwaarde biedt voor Skedify. Het onderzoek is dus in grote delen opgesplitst: Skedify, EventSourcing, huidige manier van werken en de conclusies.

In het deel over Skedify zal er onderzocht worden wat ze allemaal doen, en hoe ze dit doen. Er zal ook een gesprek komen met de business kant van Skedify om vragen te kunnen beantwoorden en om duidelijke conclusies te kunnen trekken.

In het deel over EventSourcing gaan we iets dieper in hoe EventSourcing juist werkt en hoe het gebruikt kan worden.
\section{Literatuurstudie}
\label{sec:literatuurstudie}
\section{Proof of concept}
\label{sec:proof-of-concept}

Voor de proof-of-concept zullen bepaalde functies van Skedify geschreven worden in het huidige systeem en in een EventSourcing systeem. Daarbij zal er gekeken worden naar voordelen en nadelen van beide systemen.

%% Voeg hier je eigen hoofdstukken toe die de ``corpus'' van je bachelorproef
%% vormen. De structuur en titels hangen af van je eigen onderzoek. Je kan bv.
%% elke fase in je onderzoek in een apart hoofdstuk bespreken.

%%=============================================================================
%% Skedify
%%=============================================================================

\chapter{Skedify}
\label{ch:skedify}

Skedify is een Gentse start up gespecialiseerd in het maken van afspraken. Deze bachelorproef gaat op zoek naar een antwoord of EventSourcing een meerwaarde kan bieden voor Skedify. Skedify zorgt er voor dat een persoon in eender welke sector een afspraak kan maken op zeer korte tijd voor eender welk probleem, en dit met de juiste contact persoon van het bedrijf. Skedify is een business-to-business tool om dynamisch formulieren te gaan opbouwen, in meerdere talen, voor verschillende onderwerpen. Op deze manier is de tool heel dynamisch en kan ze dus voor heel veel doeleinden worden gebruikt. Alles wordt in een relationele databank opgeslagen.

Een van de mogelijkheden die Skedify bied is het wijzigen van een afspraak, het kan dus interessant zijn om te weten hoeveel keer dit gebeurt en of er specifieke maatregelen kunnen genomen worden.
%%=============================================================================
%% CQRS
%%=============================================================================

\chapter{CQRS}
\label{ch:CQRS}

CQRS is een term uitgevonden door Greg Young en is de opvolger van CQS. CQS betekend Command Query Seperation en is een principe die uitgevonden is door Bertrand Meyer (1988). Als er gekeken wordt naar CQS is dit al een eerste vorm van goede code schrijven. CQS zorgt er voor dat getters en setters gescheiden zijn. Getters zijn strict bedoeld om een waarde uit de huidige state te halen en deze terug te geven. Setters zijn bedoeld om een wijziging te doen (of een algemene actie uit te voeren), setters geven geen waarde terug maar void.

Het grootste probleem met CQS is dat developers acties uitvoeren in getters wat voor bugs kan zorgen, bijvoorbeeld het huidige project is een spel en de totale score wordt opgevraagd. Dan ziet de getter naam er als volgt uit getTotalScore, maar als deze getter een andere methode oproept zoals calculateTotalScore, en deze methode wijzigt de state, dan kan elke keer dat getTotalScore opgeroepen wordt het resultaat anders zijn.

Een tweede probleem is dat getters en setters in een en dezelfde klasse zijn gedefinieerd. Er is geen stricte scheiding tussen de lees kant en de schrijf kant.

Dit is waar CQRS komt kijken, Command Query Responsibility Segregation. CQRS zorgt er voor dat de lees kant en de schrijf kant strict gescheiden zijn. Het zijn bijna 2 applicaties die naast elkaar staan. Een command ligt aan de schrijfkant en een query aan de leeskant. De leeskant gaat zijn informatie halen bij de databank (of een andere vorm van opslagmechanisme), dit kan via sql queries, ORM tools, enzovoort. De manier waarop dit gebeurt staat volledig los van hoe de schrijfkant communiceert met het opslagmechanisme.

Heel veel applicaties, onder andere ook die van Skedify zijn veel intensiever aan de leeskant dan aan de schrijfkant. Vermits de lees- en schrijfkant nu strict gescheiden zijn kan er gebruik gemaakt worden van schalingsmechanismen. Beter nog, de leeskant en schrijfkant kunnen individueel geschaald worden.

CQRS speelt en grote rol bij EventSourcing, vandaar dit korte hoofdstuk.
%%=============================================================================
%% EventSourcing
%%=============================================================================

\chapter{EventSourcing}
\label{ch:eventsourcing}

EventSourcing is geen nieuwe uitvinding, maar wordt al jaren gebruikt in andere sectoren. Bekende voorbeelden zijn de bankindustrie, wetgeving, en patiënten fisches die opgeslagen zijn bij dokters.

EventSourcing, is zoals de naam al verklapt, dat events de bron zijn van heel de applicatie. Het is ook zo dat de events de "single source of truth" zijn van de applicatie om andere state uit af te leiden. Een log bestaat ook uit een geschiedenis van events, maar deze wordt niet gebruikt als "single source of truth" op de huidige state af te leiden.

Als er naar de bankindustrie wordt gekeken dan is het bedrag op iemand zijn rekening niet een getal dat opgeslagen is in een databank, of toch niet de enige vorm hier van. Er is een lijst van transacties die tot dit getal komen. Het zal ook wel in een databank opgeslagen zijn, maar dit is den puur een caching mechanisme.

Het is zo dat nu pas, de laatste jaren, EventSourcing opkomt in de informatica sector. In de volgende hoofdstukken zal er dieper ingegaan worden op onderdelen van EventSourcing.
  %%=============================================================================
%% Aggregates
%%=============================================================================

\section{Aggregates}
\label{sec:aggregates}

Aggregaten zijn de entiteiten die behoren tot het domein. Een aggregaat is een entiteit die zal zorgen dat invariants (besproken in \ref{sec:invariants}) gegarandeerd zullen worden. Eens dit het geval is zal die er voor zorgen dat de juiste domain events opgeslagen worden en de state van de applicatie wordt gewijzigd.

Deze aggregaten moeten correct gekozen worden, zodat er geen gevaar is dat er 'god' aggregaten ontstaan. Een 'god' aggregaat is een aggregaat die alles in een applicatie doet, in de meeste gevallen is dit een 'User' aggregaat.

In het geval van Skedify zijn goede aggregaten: Appointment, Subject, Question, Answer, Office, ...
  %%=============================================================================
%% Value Objects
%%=============================================================================

\section{Value Objects}
\label{sec:value-objects}

Value objects zijn objecten die zorgen dat een bepaalde waarde altijd in een geldige staat is. Bijvoorbeeld, er kan een value object zijn voor een Email. Dit email adres moet ten alle tijden geldig zijn. Dit kan afgedwongen worden door in de constructor van een klasse een waarde te ontvangen en deze te controleren. Indien deze waarde fout is wordt er een exceptie gegooid. Deze exceptie zal dan opgevangen worden in de lagen daarboven.

Een belangrijke regel bij Value objects is dat deze geen identiteit bevatten, ze hebben geen unieke identifier en zijn daarom geen entiteit. Value objects kunnen wel tot een entiteit behoren zoals Bijvoorbeeld een User entiteit.

  %%=============================================================================
%% Messages
%%=============================================================================

\section{Messages}
\label{sec:messages}

Communicatie in een applicatie draait rond messages. Er zijn 3 soorten messages die terug komen in een applicatie: informational, interrogatory en imperative. \textcite{Verraes2015Messages}
Binnen een applicatie met EventSourcing wordt er ook gebruik gemaakt van messages.


\subsection{Imperative Messages}
\label{subsec:imperative-messages}

De imperative of imperatieve messages, zijn messages die de intentie hebben om de ontvanger van dit bericht een actie te laten uitvoeren.
Binnen EventSourcing zijn is een typisch voorbeeld een Command, waarbij het Command de intentie heeft om een actie uit te voeren. Binnen een systeem zijn deze Commands iets typisch dat uit de business kant komt. Wanneer het volgende voorbeeld door iemand van de business gelezen wordt, weet die ook meteen wat er zal gebeuren, als er gewerkt wordt met deze Command.

\codefragment{source/Messages-ChangeAppointmentStartDate.java}{Een voorbeeld van een Command}

Wanneer er met deze klasse gewerkt wordt (of met objecten van deze klasse), dan verwachten we dat de ontvanger iets zal wijzigen.

\subsection{Interrogatory Messages}
\label{subsec:interrogatory-messages}

De interrogatory of ondervragende messages, zijn messages die iets vragen over de huidige state van de applicatie. In elke applicatie worden er dingen opgevraagd. Binnen EventSourcing wordt dit ook gedaan aan de hand van Queries, deze Queries gaan informatie opvragen van de huidige state. Deze queries worden ook opgelegd door de business kant, de business kan, wanneer ze het voorbeeld zien, meteen begrijpen wat er zal gebeuren als er met deze Query wordt gewerkt.

\codefragment{source/Messages-MyAppointmentsForToday.java}{Een voorbeeld van een Query}

Bij dit voorbeeld is het duidelijk det er informatie opgevraagd wordt en dat er geen wijzigingen aan de huidige state zullen gebeuren.

\subsection{Informational Messages}
\label{subsec:informational-messages}

De informational of de informatieve messages, zijn messages die iets over zichzelf willen vertellen. Bij EventSourcing gaan we dit uitdrukken in DomainEvents. DomainEvents zijn de messages die opgeslagen zullen worden in de databank. DomainEvents worden ook meestal in de verleden tijd geschreven omdat ze benadrukken dat er iets gebeurt is.

\codefragment{source/Messages-AppointmentsStartDateHasBeenChanged.java}{Een voorbeeld van een DomainEvent}

Dit voorbeeld lijkt heel sterk om het voorbeeld van de Command, maar de intentie van beide klassen is totaal verschillend. De klasnamen zijn altijd zeer expliciet, soms zijn ze wat lang maar het is duidelijk wat de bedoeling is. Omdat deze zo expliciet geschreven zijn, kan er gemakkelijk mee gecommuniceerd worden wat het hele doel is van messages.

  %%=============================================================================
%% Domain Events
%%=============================================================================

\chapter{Domain Events}
\label{ch:domain-events}

Hi!
  %%=============================================================================
%% Invariants
%%=============================================================================

\chapter{Invariants}
\label{ch:invariants}

Hi!
  %%=============================================================================
%% Eventual Consistency
%%=============================================================================

\section{Eventual Consistency}
\label{sec:eventual-consistency}

...TODO...
  %%=============================================================================
%% Event Store
%%=============================================================================

\section{Event Store}
\label{sec:event-store}

De EventStore is een belangrijk concept bij EventSourcing. De EventStore is de plaats waar alle DomainEvents opgeslagen zullen worden. De EventStore kan elk soort databank zijn, er is ook een speciaal gemaakt en is te vinden op https://geteventstore.com. Voor deze bachelorproef, en de proof-of-concept zal er gebruik gemaakt worden van een simpele EventStore in mysql. Een EventStore is helemaal niet zo moeilijk qua structuur, het bevat volgende minimale velden:

\begin{itemize}
  \item{id, een simpele id, die autoincrementeel kan zijn}
  \item{stream_id, een id die bij een Aggregate hoort, dit is een GUID. Alle events die horen tot een bepaalde aggregate zullen dezelfde stream_id hebben.}
  \item{stream_version, een versie die incrementeel is per aggregaat. Dit zorgt er voor dat event in de juiste volgorde kunnen worden afgespeeld indien nodig.}
  \item{event_name, de naam van het DomainEvent dat opgeslagen moet worden. Via deze naam kan het juiste object weer opgebouwd worden in onze code.}
  \item{payload, dit bevat de effectieve data van het DomainEvent. Dit kan een json object zijn. Dit wordt mee gebruikt, naast de naam om het object weer op te bouwen.}
  \item{recorded_at, een timestamp waarop het event opgeslagen is geweest. Dit kan handig zijn voor de audit log.}
\end{itemize}

  %%=============================================================================
%% Audit Log
%%=============================================================================

\section{Audit Log}
\label{sec:audit-log}

Hi!
  %%=============================================================================
%% Projections
%%=============================================================================

\section{Projections}
\label{sec:projections}

Projections zijn een projectie van de huidige staat die berekend is van alle voorbijgaande events. Bij een systeem waar er met een relationele databank wordt gewerkt, is de databank zelf de projectie. Met als grote verschil dat daar ook de databank de enige echte bron van waarheid is. Dat is niet het geval bij een systeem met EventSourcing, daar is de bron van waarheid de events die opgeslagen zijn.

Projections hebben een heel groot voordeel dat er meerdere kunnen zijn. Er kan een databank gegenereerd worden voor de business kant die dan allerlei statistieken kan bekijken. Er kan ook een databank gegenereerd worden voor werknemers die met een bepaalde applicatie werken. Er kan zelf een databank of ander opslagsysteem gevuld worden los van de huidige databanken voor de andere partijen. Op deze manier staan we los van de databank.

Deze projecties zijn een berekening van de events, dit wilt ook zeggen dat we niet moeten optimaliseren voor de schrijfkant (zie cqrs). Er kan geoptimaliseerd worden voor de leeskant, er is geen nood om een aantal tabellen elke keer te gaan joinen, als er rechtstreeks naar 1 tabel geschrijven wordt.

Om een wijziging te kunnen doen in een aggregate, moet eerst de huidige state opgebouwd worden vanuit de DomainEvents. Dit kan echter zeer inefficiënt zijn wanneer er duizenden events beschikbaar zijn. Een projectie kan dan perfect dienen als een cache, om dit probleem op te lossen.
%%=============================================================================
%% Schaalbaarheid
%%=============================================================================

\section{Schaalbaarheid}
\label{sec:schaalbaarheid}

Een systeem dat gebruik maakt van EventSourcing, kan geschaald worden, het schalen is relatief gemakkelijk door volgende redenen.
De EventStore is een append-only database, er kan altijd een kopie genomen worden van deze database. Er moet niet gekeken worden of er al dan niet wijzigingen zijn in de data en welke de correcte data is. Synchroniseren van meerdere EventStores is een kwestie van onbrekende rijen aan te vullen. Hiervoor kan er een mechanisme gebruikt worden zoals "geef alle events na id X".
Schalen van de applicatie zelf kan als volgt gebeuren. Er kunnen projecties bijgemaakt worden die zichzelf gaan 'voeden' met alle events uit de EventStore. Van zodra ze up-to-date zijn met de EventStore kunnen ze als projection listener geregistreerd worden, om live naar de events te luisteren.
%%=============================================================================
%% Voordelen
%%=============================================================================

\section{Voordelen}
\label{sec:voordelen}

Er zijn heel wat voordelen verbonden aan EventSourcing. Een groot voordeel is dat je van technologie kan wisselen voor de projecties. Stel er werd gebruik gemaakt van MySQL om de huidige staat in op te slaan. Wanneer er gemerkt wordt dat er nood is aan een graph database in plaats van een relationele database dan kan er een projectie bijgemaakt worden die simultaan met de mysql database loopt. Zodra de graph database helemaal gevoed is met de events uit de EventStore, dan kan de MySQL databank weggegooid worden. Vermits een projectie opgebouwd is in code, en in een versiebeheersysteem zit, moet er niet nagedacht worden om een eventuele back-up te nemen van deze MySQL databank. Indien er later opnieuw van deze databank gebruik gemaakt moet worden, dan kan deze projectie terug ingeladen worden en helemaal opnieuw opgebouwd worden.
Er werd net gesproken over de graph database, stel dat de applicatie al 3 jaar draait en nu komt deze graph database er bij. Zodra de events allemaal gevoed zijn in de graph database, dan staat deze op het punt alsof ze al van in begin geprogrammeerd was. Dit is het grote voordeel van al de events bij te houden.
Een ander voordeel van EventSourcing is dat je enkel de EventStore moet backuppen. Indien elke server corrupt of kapot gaat en alle MySQL instanties zijn kapot, dan kunnen deze opnieuw opgebouwd worden. Dit is allemaal mogelijk omdat de EventStore, de single source of truth is.
% Problemen
% Kosten
% Wat wijzigen
%%=============================================================================
%% Conclusie
%%=============================================================================

\chapter{Conclusie}
\label{ch:conclusie}

%% TODO: Trek een duidelijke conclusie, in de vorm van een antwoord op de
%% onderzoeksvra(a)g(en). Wat was jouw bijdrage aan het onderzoeksdomein en
%% hoe biedt dit meerwaarde aan het vakgebied/doelgroep? Reflecteer kritisch
%% over het resultaat. Had je deze uitkomst verwacht? Zijn er zaken die nog
%% niet duidelijk zijn? Heeft het ondezoek geleid tot nieuwe vragen die
%% uitnodigen tot verder onderzoek?

\lipsum[76-80]



%%---------- Back matter ------------------------------------------------------

\phantomsection
\printbibliography

\addcontentsline{toc}{chapter}{\textcolor{maincolor}{\IfLanguageName{dutch}{Bibliografie}{Bibliography}}}

\listoffigures
\listoftables

\end{document}
